\begin{table}[H]
\centering
\caption{OP: sample construction by sector and year}
\GOneLabel{tab:deep-op-sample}
\begingroup\fontsize{9.2}{11.2}\selectfont
\setlength{\tabcolsep}{4.3pt}\renewcommand{\arraystretch}{1.1}
\begin{tabular}{llrrrrr}
\toprule
Code & Year & Source & \makecell{Eligible\\inputs} & \makecell{Consecutive\\pairs} & \makecell{Accepted\\pairs} & Residuals \\
\midrule
10 & 2011 & 567 & 75 & 0 & 0 & 75 \\
10 & 2012 & 589 & 76 & 24 & 24 & 76 \\
10 & 2013 & 597 & 79 & 20 & 20 & 79 \\
10 & 2014 & 606 & 92 & 25 & 25 & 92 \\
10 & 2015 & 628 & 0 & 0 & 0 & 0 \\
10 & 2018 & 696 & 75 & 0 & 0 & 75 \\
10 & 2019 & 666 & 82 & 16 & 16 & 82 \\
10 & 2020 & 628 & 59 & 19 & 19 & 59 \\
10 & 2021 & 730 & 51 & 11 & 11 & 51 \\
10 & 2022 & 765 & 94 & 13 & 13 & 94 \\
\bottomrule
\end{tabular}
\endgroup
\appnotes{Source counts include records before input restrictions. Eligible inputs requires the positive measurements for this method. Consecutive pairs additionally requires an eligible record in the immediately preceding year. Accepted pairs applies the fixed sector sample requirements after pooling these calendar pairs across years. Residuals counts eligible records in a sector with estimated coefficients, including records outside the dynamic stage. We use one set of sector coefficients and one pooled dynamic-sample reference mean across all years. OP requires a separate sample for its continuation calculation.}{G0.5 v3 modified manufacturing records with fictional company links; authors' calculations.}
\end{table}
